\clearpage\SourceLeaf{146}%
\pdfbookmark[1]{Gauss to von Lindenau}{lindenau1808nov}
\fancyhead[L]{\small GAUSS TO VON LINDENAU}
\begin{center}
{\small GAUSS TO VON LINDENAU.\par}
\medskip\hrule\medskip
{\small Monatliche Correspondenz zur Beförderung der Erd- und Himmels-Kunde,\par
edited by Baron \textsc{von Zach}. January 1809.\par}
\medskip\hrule
\end{center}\medskip
\begin{flushright}Göttingen, 30 November 1808.\end{flushright}
My warmest thanks for the observations of Ceres and Juno that you were so kind as to send me. I am also grateful for the comet observations from Marseille, although some \emph{very} large errors seem to have been made in them, so it will be hard to draw any conclusions from them. Today I would like to send you a few more remarks on the problem I discussed in the paper you included in the October issue of the \emph{Monatl. Corresp.} For \emph{traveling} observers, the problem I would especially recommend is different from \emph{this} one. I mentioned it at the beginning of that paper and discussed it in a programme printed here recently, a copy of which I am pleased to enclose. The problem treated in the \emph{Monatl. Corresp.} will be particularly useful to observers who want to determine the position of their observing site as precisely as possible with an instrument that is less than perfect, and at the same time determine the errors of the instrument. In that case it is quite easy to find plenty of suitable stars for any hour if one has calculated a table of altitudes in advance for the most readily recognizable stars. The table may be useful in other ways too and can very conveniently be put on a chart. There is no need to restrict oneself to stars of the first magnitude.
\SourceLeaf{147}%
You can see from my example that stars of the second magnitude can also be used. With some practice, if there is no moonlight, one can even observe stars of the third and fourth magnitudes readily enough, although of course more care is then needed to avoid \emph{mistaking one star for another}. As for my solution itself, it is the one I hit on at once on my own when I looked for a method as convenient as possible to use in practice. It did not occur to me to look elsewhere for solutions to a problem that, as far as I knew, no one had tackled before. Only after I had sent off my paper did I notice that another problem, though very different in purpose, is essentially identical as far as the solution goes. It is the problem of finding the orientation of the Sun’s equator and at the same time the declination of a sunspot from three \emph{heliocentric} positions of the spot. The latter quantity corresponds to the altitude $h$ in our problem. Many geometers, as is well known, have tackled this problem, but none has given so elegant a solution as \textsc{Cagnoli}. \textsc{La Lande} gives it in his \textit{Astronomie}. For elegance I prefer this solution to my own, although I believe the solution itself can be derived from the three fundamental equations more elegantly than \textsc{Cagnoli} derived it. For ease of practical use, too, all the other solutions that \textsc{La Lande} gives fall well short of \textsc{Cagnoli}’s. Mine would be a little more convenient if one wanted only $\varphi$ and $k$, while \textsc{Cagnoli}’s would be a little shorter if one wanted $h$ as well. In the first case I need 18 logarithms and \textsc{Cagnoli} 21. In the second I need 24 and \textsc{Cagnoli} again only 21. It is worth noting, though, that in \textsc{Cagnoli}’s solution all 21 refer to almost as many angles, namely 19 \emph{different} angles. In my solution, by contrast, for many angles one looks up the sine and cosine, or other trigonometric functions, together. (Thus one looks at 14 different places in the tables if only $\varphi$ and $k$ are wanted and at 18 if $h$ is wanted as well.) That certainly makes a difference, especially when using \textsc{Taylor}’s excellent tables, with which almost the whole job comes down to looking things up. The formulas for our problem according to \textsc{Cagnoli}’s solution were as follows.

Compute the three auxiliary angles $A,A',A''$ from the equations
\SourceLeaf{148}%
\[
\begin{aligned}
\operatorname{tang}A&=\frac{\sin\frac12(\delta''-\delta')}{\cos\frac12(\delta''+\delta')}\operatorname{cotang}\tfrac12(t''-t')\\
\operatorname{tang}A'&=\frac{\sin\frac12(\delta-\delta'')}{\cos\frac12(\delta+\delta'')}\operatorname{cotang}\tfrac12(t-t'')\\
\operatorname{tang}A''&=\frac{\sin\frac12(\delta'-\delta)}{\cos\frac12(\delta'+\delta)}\operatorname{cotang}\tfrac12(t'-t)
\end{aligned}
\]
then set
\[
\operatorname{tang}B=\frac{\sin\frac12(\delta'-\delta)}{\cos\frac12(\delta'+\delta)}\operatorname{cotang}(A-A')
\]
which gives
\[
k=\tfrac12(t'+t)-B
\]
Next set $A'+A''-A=C$ and take
\[
\begin{aligned}
\operatorname{tang}D&=\frac{\cos\frac12(t-k+C)}{\cos\frac12(t-k-C)}\operatorname{tang}(45^0+\tfrac12\delta)\\
\operatorname{tang}E&=\frac{\sin\frac12(t-k-C)}{\sin\frac12(t-k+C)}\operatorname{cotg}(45^0+\tfrac12\delta)
\end{aligned}
\]
\expandafter\def\csname HistoricalNoteText6\endcsname{[Handwritten remark:]
\[
\begin{aligned}
\frac{\cos\varphi}{\cos h}\sin(t-k)&=\sin(A'+A''-A)\\
\frac{\cos\varphi}{\cos h}\sin(t'-k)&=\sin(A''+A-A')\\
\frac{\cos\varphi}{\cos h}\sin(t''-k)&=\sin(A+A'-A'')
\end{aligned}
\]}
which gives\HistoricalNote{6}
\[
\varphi=D+E,\qquad h=D-E
\]

One inconvenience of this method is that it is hard to give a definite rule in advance for which semicircle to take the angles $A,A',A'',B,D,E$ in, if one calculates from the analytical formulas alone without a diagram and uses only the clock readings from the observations, so that one does not know in advance on which sides of the meridian the three stars were observed. As is well known, the tangents alone leave the choice undecided. It can be shown, however, that the choice can be made arbitrarily at first. In some cases one will then have to make the following changes.

1) Replace $k$ by $k+180^0$ or, which amounts to the same thing here, $k-180^0$, if the values obtained for $\varphi$ and $h$ give $\cos\varphi$ and $\sin h$ opposite signs.
\SourceLeaf{149}%

2) If the values found for $h$ and $\varphi$ are outside the limits $0$ and $90^0$, use their differences from the nearest multiple of $180^0$.

3) Take the elevation of the pole as north or south according as $\sin\varphi$ and $\sin h$ have the same or opposite signs.

For instance, if the formulas above gave $\varphi=231^0$, $h=-127^0$, the actual elevation of the pole would be $51^0$ and the actual altitude $53^0$. The value of $k$ would stay as it was. If instead one found $\varphi=231^0$, $h=+127^0$, one would have to increase or decrease $k$ by $180^0$. This point is of course mentioned solely for analytical completeness. If $k$ needs to be changed by $180^0$, it will already be clear, since the clock offset is never uncertain by as much as 12 hours.

For comparison I also enclose the numerical calculation for the same example using the formulas above.

\clearpage\begin{center}\emph{Calculation of the Example by Cagnoli’s Method.}\end{center}
\begin{center}\small
\begin{minipage}[t]{.49\textwidth}\centering
\begin{tabular}[t]{@{}lr@{}}
$\tfrac12(\delta''-\delta')=$&$-24^0\ 49'\ 59''55$\\
$\tfrac12(\delta-\delta'')=$&$-5\quad17\quad25.90$\\
$\tfrac12(\delta'-\delta)=$&$+30\quad7\quad25.45$\\[10pt]
$\tfrac12(t''-t')=$&$-129\quad41\quad39.45$\\
$\tfrac12(t-t'')=$&$+135\quad0\quad4.72$\\
$\tfrac12(t'-t)=$&$-5\quad18\quad25.27$
\end{tabular}
\end{minipage}\hfill
\begin{minipage}[t]{.49\textwidth}\centering
\begin{tabular}[t]{@{}lr@{}}
$\tfrac12(\delta''+\delta')=$&$63^0\ 27'\ 6''15$\\
$\tfrac12(\delta+\delta'')=$&$33\quad19\quad40.70$\\
$\tfrac12(\delta'+\delta)=$&$58\quad9\quad40.25$
\end{tabular}
\end{minipage}
\par\medskip
\begin{tabular}[b]{@{}lr@{}}
$\log\sin\tfrac12(\delta''-\delta')\;.\;.\;.$&$9.6232267\,\mathrm n$\\
$\mathrm{Compl.}\log\cos\tfrac12(\delta''+\delta')\;.$&$0.3497391$\\
$\log\operatorname{cotang}\tfrac12(t''-t)\;.\;.\;.$&$9.9191030$\\\hline
$\log\operatorname{tang}A\;.\;.\;.\;.\;.\;.$&$9.8920688\,\mathrm n$
\end{tabular}
\qquad $A=142^0\quad2\quad50''70$
\par\medskip
\begin{tabular}[b]{@{}lr@{}}
$\log\sin\tfrac12(\delta-\delta'')\;.\;.\;.\;.$&$8.9647590\,\mathrm n$\\
$\mathrm{Compl.}\log\cos\tfrac12(\delta+\delta'')\;.$&$0.0780332$\\
$\log\operatorname{cotang}\tfrac12(t-t')\;.\;.\;.$&$0.0000199\,\mathrm n$\\\hline
$\log\operatorname{tang}A'\;.\;.\;.\;.\;.$&$9.0428121$
\end{tabular}
\qquad $A'=6^0\ 17'\ 51''34$
\end{center}
\SourceLeaf{150}%
\begin{center}\begingroup\small\setlength{\tabcolsep}{2.5pt}
\sbox{0}{\begin{tabular}{@{}lr@{\hspace{16pt}}rl@{}}
$\log\sin\tfrac12(\delta'-\delta)\;.\;.\;.\;.$&$9.7005892$&&\\
$\mathrm{Compl.}\log\cos\tfrac12(\delta'+\delta)\;.$&$0.2777516$&&\\\cline{1-2}
$(*)\;.\;.\;.\;.\;.\;.\;.\;.$&$9.9783408$&&\\
$\log\operatorname{cotang}\tfrac12(t'-t)\;.\;.\;.$&$1.0320274$&&\\\cline{1-2}
$\log\operatorname{tang}A''\;.\;.\;.\;.\;.$&$1.0103682\,\mathrm n$&$A''=$&$95^0\ 34'\ 36''24$\\
&&$A-A'=$&$135\quad44\quad59.36$\\
&&$C=$&$-40\quad10\quad23.12$\\[7pt]
$(*)\;.\;.\;.\;.\;.\;.\;.\;.$&$9.9783408$&&\\
$\log\operatorname{cotang}(A-A')\;.\;.\;.$&$0.0113684\,\mathrm n$&&\\\cline{1-2}
$\log\operatorname{tang}B\;.\;.\;.\;.\;.$&$9.9897092\,\mathrm n$&$B=$&$315^0\ 40'\ 43''55$\\
&&$\tfrac12(t+t')=$&$318\quad24\quad44.78$\\\cline{3-4}
&&$k=$&$2\quad44\quad1.23$\\[4pt]
&&$\tfrac12(t-k+C)=$&$140\quad24\quad22.85$\\
&&$\tfrac12(t-k-C)=$&$180\quad34\quad45.97$\\
&&$45^0+\tfrac12\delta=$&$59\quad1\quad7.40$\\[7pt]
$\log\cos\tfrac12(t-k+C)\;.\;.$&$9.8868199\,\mathrm n$&&\\
$\mathrm{Compl.}\log\cos\tfrac12(t-k-C)$&$0.0000222\,\mathrm n$&&\\
$\log\operatorname{tang}(45^0+\tfrac12\delta)\;.\;.\;.$&$0.2215478$&&\\\cline{1-2}
$\log\operatorname{tang}D\;.\;.\;.\;.\;.$&$0.1083899$&$D=$&$52^0\ 4'\ 36''35$\\[7pt]
$\log\sin\tfrac12(t-k-C)\;.\;.$&$8.0048736\,\mathrm n$&&\\
$\mathrm{Compl.}\log\sin\tfrac12(t-k+C)$&$0.1956297$&&\\
$\log\operatorname{cotang}(45^0+\tfrac12\delta)\;.\;.$&$9.7784522\,\mathrm n$&&\\\cline{1-2}
$\log\operatorname{tang}E\;.\;.\;.\;.\;.$&$7.9789555\,\mathrm n$&$E=$&$-0^0\ 32'\ 45''02$\\
&&Thus $\varphi=$&$51\quad31\quad51.33$\\
&&$h=$&$52\quad37\quad21.37$\\
\end{tabular}}
\ifdim\wd0>\linewidth\resizebox{\linewidth}{!}{\usebox{0}}\else\usebox{0}\fi
\endgroup\end{center}
\par\bigskip\centerline{\rule{.27\linewidth}{.4pt}}
